\documentclass[../DoAn.tex]{subfiles}
\begin{document}
\section{Khảo sát hiện trạng}
\label{section:1.1}
Hiện nay, khác với các nước đã có nền giáo dục số phát triển và ứng dụng công nghệ AI rộng rãi trong học tập, tại Việt Nam, đa số các hệ thống học tập flashcard vẫn là các ứng dụng đơn giản với chức năng cơ bản là tạo, lưu trữ và hiển thị flashcard, ít có tính năng thông minh, chủ yếu mang tính quản lý nội dung tĩnh. Bên cạnh đó, đã có nhiều hệ thống flashcard động với giao diện, bố cục và cách thức quản lý đa dạng, ví dụ như Anki với thuật toán spaced repetition, Quizlet với tính năng học tập tương tác và chia sẻ bộ flashcard, Memrise với tích hợp đa phương tiện, hoặc các ứng dụng flashcard trong nước như StudyNow, FlashcardVN, và các hệ thống học từ vựng tiếng Anh như eJOY, Monkey Junior.

Tuy nhiên, thực tế cho thấy phần lớn các hệ thống flashcard hiện có, đặc biệt là các hệ thống quốc tế như Anki và Quizlet, mặc dù có tính năng spaced repetition và hoạt động ổn định, nhưng việc tùy biến và tích hợp AI để hỗ trợ học tập theo ngữ cảnh còn hạn chế. Các hệ thống này chủ yếu dựa trên thuật toán spaced repetition cơ bản mà chưa tích hợp trợ lý ảo AI để sinh bài test, quiz hoặc đoạn văn minh họa tự động dựa trên nội dung flashcard đang học. Ở Việt Nam, các hệ thống flashcard trong nước thường được phát triển bởi các công ty công nghệ giáo dục với chi phí ban đầu có thể cao, nhưng việc duy trì và cập nhật tính năng mới, đặc biệt là tích hợp AI, còn phụ thuộc rất nhiều vào nguồn lực và chiến lược phát triển của từng công ty. Điều này dẫn đến việc người học thường xuyên phải sử dụng nhiều công cụ khác nhau để đáp ứng các nhu cầu học tập khác nhau, gây bất tiện và giảm hiệu quả học tập.

Bên cạnh đó, nhu cầu học tập cá nhân hóa với sự hỗ trợ của AI đang là vấn đề được nhiều người học quan tâm. Các hệ thống hiện tại chưa đáp ứng tốt nhu cầu tương tác với AI trong quá trình học để đặt câu hỏi, nhận giải thích tức thì, hoặc yêu cầu sinh nội dung học tập phù hợp như bài test, quiz và đoạn văn minh họa dựa trên flashcard đang học. Đối với giáo viên, các hệ thống hiện tại cũng chưa cung cấp đầy đủ công cụ quản lý lớp học và theo dõi tiến độ học tập chi tiết của học viên, gây khó khăn trong việc đánh giá và hỗ trợ học sinh một cách hiệu quả.

Dựa trên các khảo sát và đánh giá trên, đề tài hướng tới xây dựng một nền tảng học tập flashcard với kiến trúc client-server hiện đại, tích hợp trợ lý ảo AI và thuật toán SM-2 để tối ưu hóa quá trình ghi nhớ. Hệ thống sẽ cung cấp các tính năng quản lý flashcard thông minh và hỗ trợ người học trong việc học tập cá nhân hóa, tương tác với AI để sinh nội dung học tập phù hợp, và quản lý lớp học cho giáo viên. Việc sử dụng nền tảng này không chỉ mang lại sự tiện lợi, hiệu quả mà còn giúp người học có trải nghiệm học tập thông minh và cá nhân hóa, đồng thời cung cấp công cụ quản lý và theo dõi tiến độ học tập toàn diện cho cả người học và giáo viên.

\section{Đặt vấn đề}
\label{section:1.2}
Trong bối cảnh chuyển đổi số mạnh mẽ, các mô hình học tập truyền thống đang dần được thay thế bởi các nền tảng học trực tuyến nhằm đáp ứng nhu cầu học tập linh hoạt, cá nhân hóa và không giới hạn không gian, thời gian. Tuy nhiên, phần lớn các hệ thống học tập flashcard hiện nay vẫn chỉ dừng lại ở mức quản lý và hiển thị nội dung, chưa thực sự hỗ trợ người học trong quá trình hiểu – ghi nhớ – ôn luyện – đánh giá năng lực một cách thông minh và hiệu quả. Các hệ thống flashcard truyền thống không có thuật toán tự động điều chỉnh khoảng thời gian ôn tập dựa trên khả năng nhớ lại của từng cá nhân, dẫn đến việc người học phải ôn tập theo lịch cố định không phù hợp với khả năng ghi nhớ thực tế, gây lãng phí thời gian vào những nội dung đã thuộc hoặc bỏ sót những kiến thức cần được củng cố kịp thời. Bên cạnh đó, người học không thể đặt câu hỏi hoặc nhận giải thích ngay lập tức khi gặp khó khăn với một khái niệm trong quá trình học, làm giảm hiệu quả học tập và tăng cảm giác cô đơn trong quá trình tự học. Hệ thống cũng không thể tự động sinh bài test, quiz hoặc đoạn văn minh họa phù hợp với nội dung flashcard đang học, khiến người học thiếu công cụ để củng cố và mở rộng kiến thức một cách có hệ thống. Hơn nữa, nội dung học tập thiếu tính thích ứng với năng lực và tiến trình riêng của từng cá nhân do hệ thống không theo dõi và phân tích hiệu suất học tập để đề xuất flashcard cần ôn lại hoặc điều chỉnh lịch học phù hợp. Đối với giáo viên và người học, việc theo dõi tiến độ học tập chi tiết, phân tích điểm mạnh và điểm yếu, cũng như quản lý lớp học một cách hiệu quả vẫn còn nhiều hạn chế do thiếu công cụ thống kê và phân tích chuyên sâu.

Sự phát triển của trí tuệ nhân tạo (AI), đặc biệt là các mô hình ngôn ngữ lớn (LLM) và các nền tảng workflow AI như Dify, đã mở ra cơ hội xây dựng nền tảng học tập flashcard thông minh, nơi người học được hỗ trợ bởi Trợ lý Ảo (AI Tutor) có khả năng hiểu, phản hồi và hướng dẫn phù hợp với năng lực học tập của từng cá nhân, đồng thời tích hợp thuật toán spaced repetition như SM-2 (SuperMemo 2) để tự động điều chỉnh khoảng thời gian ôn tập dựa trên Easiness Factor và recall score của từng flashcard, tối ưu hóa quá trình ghi nhớ. Vì vậy, việc phát triển hệ thống Web Flashcard tích hợp AI Tutor thông qua Dify workflows và thuật toán SM-2 là cần thiết để khắc phục các hạn chế trên, đồng thời nâng cao hiệu quả và chất lượng học tập trong môi trường trực tuyến thông qua việc cá nhân hóa quá trình học tập và cung cấp công cụ hỗ trợ thông minh cho cả người học và giáo viên.

\section{Mục tiêu và phạm vi đề tài}
\label{section:1.3}
Mục tiêu chính của đề tài là xây dựng một hệ thống web học tập flashcard tích hợp trợ lý ảo AI, cung cấp nền tảng học tập thông minh và cá nhân hóa cho ba vai trò người dùng chính là Student, Teacher và Admin. Hệ thống nhằm tích hợp trợ lý ảo AI hỗ trợ học tập theo ngữ cảnh, cho phép người học tương tác với AI trong quá trình học để đặt câu hỏi, yêu cầu sinh bài test, quiz và đoạn văn minh họa dựa trên nội dung flashcard đang học. Đề tài hướng đến việc đề xuất nội dung học dựa trên hiệu suất và tiến độ người dùng thông qua việc phân tích dữ liệu học tập và tự động đề xuất flashcard cần ôn lại. Đối với giáo viên, hệ thống cung cấp công cụ quản lý lớp học và thống kê học tập chi tiết, cho phép theo dõi tiến độ học tập của toàn bộ lớp. Cuối cùng, đề tài đặt mục tiêu thiết kế kiến trúc backend mở rộng và dễ phát triển lâu dài.

Phạm vi của đề tài bao gồm việc phát triển một hệ thống web học tập flashcard với ba vai trò người dùng chính. Đối với Student, hệ thống cung cấp các chức năng quản lý flashcard cá nhân, học tập với thuật toán SM-2 để tối ưu hóa spaced repetition, tương tác với trợ lý ảo AI để sinh bài test, quiz và đoạn văn minh họa, xem thống kê tiến độ học tập cá nhân, và tham gia lớp học. Đối với Teacher, hệ thống cung cấp toàn bộ chức năng của Student cộng thêm khả năng tạo và quản lý lớp học, quản lý danh sách học viên, và xem thống kê tiến độ học tập của toàn bộ lớp. Đối với Admin, hệ thống cung cấp chức năng quản trị toàn hệ thống bao gồm xem thống kê tổng thể, theo dõi số lượng người dùng, hoạt động học tập và hiệu suất sử dụng hệ thống. Về mặt kỹ thuật, phạm vi bao gồm việc tích hợp Dify workflows để xử lý các yêu cầu AI, triển khai thuật toán SM-2 để tự động điều chỉnh khoảng thời gian ôn tập, xây dựng hệ thống theo dõi tiến độ học tập, và phát triển hệ thống quản lý test. Hệ thống được xây dựng với backend sử dụng FastAPI và Python, frontend sử dụng React Router v7 và TypeScript, database PostgreSQL với SQLModel, và xác thực qua JWT tokens.

\section{Định hướng giải pháp}
\label{section:1.4}
Định hướng giải pháp của đề tài được xây dựng dựa trên kiến trúc client-server với backend module hóa, đảm bảo sự ổn định, bảo mật và khả năng mở rộng. Backend được phát triển bằng FastAPI và Python, sử dụng SQLModel để quản lý database và Alembic cho migrations, với cấu trúc module rõ ràng bao gồm các routes cho learning, studysets, tests, chatbot, classes, categories và admin, các services chuyên biệt như LearningService cho thuật toán SM-2, GenerationService cho tích hợp AI, ChatbotService cho tương tác với AI, và AnalyticsService cho thống kê. Hệ thống tích hợp Dify workflows thông qua DifyService để xử lý các yêu cầu AI, cho phép sinh bài test, quiz và đoạn văn minh họa dựa trên nội dung flashcard, đồng thời cung cấp chatbot để trả lời câu hỏi của người học trong quá trình học. Thuật toán SM-2 được triển khai trong LearningService với việc tự động tính toán Easiness Factor, interval và next\_review\_date dựa trên recall score của từng flashcard, đảm bảo tối ưu hóa quá trình ghi nhớ. Frontend được xây dựng bằng React Router v7 và TypeScript, sử dụng Redux để quản lý state, Tailwind CSS và Chakra UI cho giao diện, đảm bảo trải nghiệm người dùng thân thiện và hiện đại. Hệ thống được thiết kế với khả năng tự host workflow AI thông qua Dify, cho phép kiểm soát logic xử lý, chi phí và hiệu năng một cách tối ưu, đồng thời đảm bảo tính bảo mật thông qua JWT authentication và OAuth2. Kiến trúc này cho phép hệ thống dễ dàng mở rộng thêm các tính năng học tập, khả năng AI và phân tích trong tương lai mà không cần thay đổi cấu trúc cốt lõi.

\section{Bố cục đồ án}
\label{section:1.5}
Phần còn lại của báo cáo đồ án tốt nghiệp này được tổ chức như sau.

Chương 2 trình bày về cơ sở lý thuyết và các công nghệ liên quan đến đề tài, bao gồm các khái niệm về spaced repetition và thuật toán SM-2, các mô hình ngôn ngữ lớn và ứng dụng trong giáo dục, kiến trúc hệ thống web hiện đại với FastAPI và React, cũng như các công nghệ hỗ trợ như PostgreSQL, JWT authentication và Dify workflows.

Chương 3 trình bày về phân tích và thiết kế hệ thống, bao gồm phân tích yêu cầu chức năng và phi chức năng, thiết kế kiến trúc tổng thể của hệ thống, thiết kế cơ sở dữ liệu với các model và quan hệ giữa chúng, thiết kế các API endpoints và services, cũng như thiết kế giao diện người dùng.

Chương 4 trình bày về triển khai hệ thống, bao gồm triển khai backend với các routes, services và tích hợp Dify workflows, triển khai thuật toán SM-2 trong LearningService, triển khai frontend với các components và pages, tích hợp authentication và authorization, cũng như các tính năng quản lý flashcard, học tập, test và chatbot.

Chương 5 trình bày về kiểm thử và đánh giá hệ thống, bao gồm các phương pháp kiểm thử đã áp dụng, kết quả kiểm thử các chức năng chính, đánh giá hiệu năng và khả năng mở rộng của hệ thống, cũng như so sánh với các hệ thống tương tự và đưa ra các hướng phát triển trong tương lai.

\end{document}

